\documentclass[11pt,a4paper]{article}
\usepackage{amsmath,amssymb}
\usepackage{listings}
\usepackage{hyperref}
\usepackage[margin=1in]{geometry}
\usepackage{fancyhdr}

\pagestyle{fancy}
\fancyhf{}
\fancyhead[L]{Module 17: Machine Learning in C}
\fancyhead[R]{C Programming Course}
\fancyfoot[C]{\thepage}

\lstset{
  language=C,
  basicstyle=\ttfamily\small,
  keywordstyle=\bfseries,
  commentstyle=\itshape,
  numbers=left,
  numberstyle=\tiny,
  frame=single,
  breaklines=true,
  tabsize=4
}

\title{Module 17: Machine Learning in C}
\author{C Programming Course}
\date{}

\begin{document}
\maketitle
\thispagestyle{fancy}

\section{Introduction to ML Concepts}
Machine learning finds patterns in data. Key paradigms: supervised learning
(labeled data), unsupervised learning (no labels), and reinforcement learning
(reward signals).

\section{Perceptron Implementation}
A perceptron is the simplest neural network: a single neuron with weighted inputs.

\begin{lstlisting}
#include <stdio.h>

double dot(const double *w, const double *x,
           int n) {
    double sum = 0.0;
    for (int i = 0; i < n; i++)
        sum += w[i] * x[i];
    return sum;
}

int predict(const double *w, const double *x,
            int n, double bias) {
    return (dot(w, x, n) + bias) > 0.0 ? 1 : 0;
}
\end{lstlisting}

\section{Activation Functions}
Common activation functions:
\begin{itemize}
  \item \textbf{Sigmoid}: $\sigma(x) = \frac{1}{1 + e^{-x}}$
  \item \textbf{ReLU}: $f(x) = \max(0, x)$
  \item \textbf{Tanh}: $\tanh(x) = \frac{e^x - e^{-x}}{e^x + e^{-x}}$
\end{itemize}

\begin{lstlisting}
#include <math.h>

double sigmoid(double x) {
    return 1.0 / (1.0 + exp(-x));
}

double relu(double x) {
    return x > 0.0 ? x : 0.0;
}
\end{lstlisting}

\section{Forward and Backward Propagation}
\textbf{Forward pass}: compute outputs layer by layer.
\textbf{Backward pass} (backpropagation): compute gradients of the loss with
respect to each weight using the chain rule:
\[
\frac{\partial L}{\partial w} =
\frac{\partial L}{\partial a} \cdot
\frac{\partial a}{\partial z} \cdot
\frac{\partial z}{\partial w}
\]

\section{Training Neural Networks}
The training loop:
\begin{enumerate}
  \item Forward pass to compute predictions.
  \item Compute the loss (e.g., mean squared error).
  \item Backward pass to compute gradients.
  \item Update weights: $w \leftarrow w - \eta \cdot \nabla L$.
\end{enumerate}

Learning rate $\eta$ controls the step size. Too large causes divergence;
too small causes slow convergence.

\end{document}
